% !TEX root = main.tex
% =============================================================================
% Cost functions
% =============================================================================

\section{Cost Functions}
\label{sec:costs}

Alg.~\ref{alg:refine} minimizes a scalar soft objective
\begin{equation}
\label{eq:J}
J(x)
=
J_{\mathrm{path}}(x)
+
J_{\mathrm{clear}}(x)
+
J_{\mathrm{union}}(x)
+
J_{\mathrm{space}}(x).
\end{equation}
Terms below are evaluated on Hermite samples (and proximity pairs from
Alg.~\ref{alg:refine}).

\subsection{Path costs}
\label{sec:costs:path}

We set
\begin{equation}
J_{\mathrm{path}}
=
\alpha\,L
+
\beta\,K.
\end{equation}
Let $\mathcal{E}$ be the set of Hermite segments (trunk and junction) and
$\mathcal{J}$ the set of unordered segment--segment joints discovered by walking
usages.

\paragraph{Aggregation (not a naive usage sum).}
Chord energy charges each physical segment once,
\begin{equation}
L
=
\sum_{e\in\mathcal{E}}
L_e,
\end{equation}
so a shared trunk is not multiplied by how many usages traverse it.
Curvature sums segment interiors once and every distinct usage joint once,
\begin{equation}
K
=
\sum_{e\in\mathcal{E}}
K_e^{\mathrm{int}}
+
\sum_{\{e,e'\}\in\mathcal{J}}
K_{e,e'}^{\mathrm{joint}}.
\end{equation}
Walking $\mathcal{U}$ is still required for $K$: the hairpin at a shared $Y$
lives at \emph{which} segments meet, even when trunk interiors were already
scored.
A naive $\sum_u L_u$ (or $\sum_u K_u$ over full usage polylines) would
over-count shared material and over-pull trunks.

\paragraph{Chord energy.}
On a sample chain $\{x_i\}$, we use squared chords rather than
$\sum_i\|x_{i+1}-x_i\|_2$:
\begin{equation}
L_e
=
\sum_{i}
\|x_{i+1}-x_i\|_2^{2}
\end{equation}
(with optional sample weights $\bar w_i=\tfrac12(w_i+w_{i+1})$).
Squaring penalizes uneven gap distributions more strongly than raw chord length
(one long chord costs more than several short ones of equal total length).

\paragraph{Radius-aware bend.}
The same centerline turn is worse for a fatter tube, so we penalize a
dimensionless bend related to $\kappa r$, not bare $\kappa$.
With unit chords $t_i$ and turning angle
\begin{equation}
\theta_i
=
\mathrm{atan2}\big(\|t_{i-1}\times t_i\|_2,\; t_{i-1}\cdot t_i\big)
\in[0,\pi]
\end{equation}
(preferred over $\arccos(t_{i-1}\cdot t_i)$, which is ill-behaved as
$\theta\!\to\!0$), the bend energy on a chain is
\begin{equation}
\label{eq:Kbend}
\sum_{i}
\frac{(\theta_i\, r_i)^{2}}{\Delta\ell_i+\varepsilon}
\approx
\int (\kappa r)^{2}\,\mathrm{d}s,
\end{equation}
with chord lengths $\Delta s_i=\|x_{i+1}-x_i\|_2$ and
$\Delta\ell_i=\tfrac12(\Delta s_{i-1}+\Delta s_i)$.
The same kernel is applied on segment interiors ($K_e^{\mathrm{int}}$) and on
three-point windows at usage joints ($K_{e,e'}^{\mathrm{joint}}$).

\subsection{Clearance}
\label{sec:costs:clear}

For two capsule / ball witnesses with centers $a,b$ and radii $r_a,r_b$, let
$d_{\mathrm{req}}=r_a+r_b+\delta_{\mathrm{margin}}$ and
\begin{equation}
\phi(a,b;r_a,r_b)
=
\Big(
\big[
d_{\mathrm{req}}-\|a-b\|_2
\big]_+^{\sigma}
\Big)^{2},
\end{equation}
where $[\cdot]_+^\sigma$ is a softplus hinge.
Then
\begin{equation}
J_{\mathrm{clear}}
=
J_{\mathrm{obs}}
+
J_{\mathrm{grp}},
\end{equation}
with $J_{\mathrm{obs}}$ summing $\phi$ over tube--obstacle pairs in
$\mathcal{N}_{\mathrm{obs}}$ and $J_{\mathrm{grp}}$ summing $\phi$ over
cross-group pairs in $\mathcal{N}_{\mathrm{grp}}$
(Alg.~\ref{alg:refine}).
Weights $\mu_{\mathrm{obs}}$ and $\mu_{\mathrm{grp}}$ scale the two sums.

\paragraph{Hard accept.}
After the loop we recompute residuals on obstacle / cross-group pairs
\emph{without} the soft margin used inside $J_{\mathrm{clear}}$ and accept only
if the hard hit count is zero (Alg.~\ref{alg:refine}).
Capsule proxies are incomplete relative to exact swept solids, but sufficient
for the geometric refine before CAD.

\subsection{Junction union}
\label{sec:costs:union}

Private junction patches are under-determined by $L$ and $K$ alone.
Inside each complex $Z$, let $\gamma_u^{Z}$ be the samples of usage $u$ on
junction edges of $Z$, and let $\mathcal{U}_Z$ be the usages that enter $Z$.
For a point $x$ with radius $r_x$ and a curve $\gamma$ with nearest-sample
radius $r_\star$, set $r_{\Sigma}=r_x+r_\star$ and
\begin{equation}
\psi(d;r_{\Sigma})
=
\begin{cases}
\displaystyle\bigl(d/r_{\Sigma}\bigr)^{2}
  & d < \alpha\, r_{\Sigma}, \\[0.4em]
0 & \text{otherwise,}
\end{cases}
\end{equation}
with cutoff factor $\alpha\sim O(1)$.
The pairwise Chamfer
\begin{equation}
\mathrm{Ch}(\gamma_u^{Z},\gamma_v^{Z})
=
\sum_{x\in\gamma_u^{Z}}
\psi\big(d(x,\gamma_v^{Z});r_{\Sigma}\big)\,\Delta s(x)
\;+\;
\sum_{y\in\gamma_v^{Z}}
\psi\big(d(y,\gamma_u^{Z});r_{\Sigma}\big)\,\Delta s(y)
\end{equation}
then defines
\begin{equation}
J_{\mathrm{union}}
=
\lambda
\sum_{Z}
\sum_{\{u,v\}\subset\mathcal{U}_Z}
\mathrm{Ch}(\gamma_u^{Z},\gamma_v^{Z}).
\end{equation}
While usages remain within about one radius sum, $\psi$ pulls them into a
compact bundle; once they separate beyond the cutoff (local split,
Fig.~\ref{fig:y-shape}b), the force releases.
Radius normalization keeps $\lambda$ comparable across tube sizes.

\subsection{Spacing regularizer}
\label{sec:costs:space}

Hermite segments are sampled in the parameter $t\in[0,1]$, not in arc length.
Without a spacing term, the optimizer can stretch $\|\gamma'(t)\|$ so that large
gaps open between samples: soft clearance and bend then see only a sparse,
biased subset of the curve.
We therefore use
\begin{equation}
J_{\mathrm{space}}
=
\eta
\sum_{i}
\big(\Delta s_i-\overline{\Delta s}\big)^{2},
\end{equation}
where $\overline{\Delta s}$ is the mean chord on the chain, to keep sample gaps
comparable and reduce this parameterization artifact.
